% TOP_PROBLEM: {"id": 3280, "title": "Decomposing k-arc-strong tournament into k spanning strong digraphs", "category_id": 3, "category": {"id": 3, "name": "graph_theory", "display_name": "Graph Theory", "description": "Problems involving graphs, networks, and their properties.", "slug": "graph-theory", "order_index": 3, "created_at": "2026-07-31T15:26:25.671Z"}, "status": "open", "problem_number": "OPG-47651", "set_id": 12}
% FIRST_POSTED: 2026-10-02
% ATTEMPT_STATUS: unresolved
% UnsolvedMath ID 3280; source catalogue code OPG-47651
% !TeX program = lualatex
% GPT-6 Astra Ultra research attempt; independently checked partial work, 2026-10-02.
\documentclass[11pt,a4paper]{article}
\usepackage[margin=25mm]{geometry}
\usepackage{fontspec}
\setmainfont{lmroman10-regular.otf}[BoldFont=lmroman10-bold.otf,
 ItalicFont=lmroman10-italic.otf,BoldItalicFont=lmroman10-bolditalic.otf]
\usepackage{amsmath,amssymb,mathtools}
\usepackage{unicode-math}
\setmathfont{latinmodern-math.otf}
\AtBeginDocument{\let\setminus\smallsetminus}
\usepackage{xurl,hyperref}
\hypersetup{colorlinks=true,urlcolor=blue}
\setcounter{secnumdepth}{0}
\setlength{\parindent}{0pt}
\setlength{\parskip}{5pt}
\setlength{\emergencystretch}{3em}
\begin{document}
\section*{Decomposing k-arc-strong tournament into k spanning strong digraphs}\label{3280}
{\small UnsolvedMath ID 3280 · OPG-47651 · Graph Theory · OpenGarden}
\subsection{Definitions and mathematical statement}
Let $T=(V,A)$ be a finite tournament: for each unordered pair of distinct
vertices exactly one orientation is an arc. A digraph is strongly
connected if a directed path joins every ordered pair of vertices.
For $k\ge1$, the tournament is $k$-arc-strong if it remains strongly
connected after deletion of any set of fewer than $k$ arcs. Equivalently,
for every nonempty proper $X\subset V$, at least $k$ arcs leave $X$.

The conjecture asks whether every $k$-arc-strong tournament admits a
partition of its arc set
\[
                      A=A_1\sqcup\cdots\sqcup A_k
\]
such that each spanning subdigraph $(V,A_i)$ is strongly connected.
Each part uses the full vertex set. These subdigraphs need not themselves
be tournaments, since most pairs of vertices may have no arc in a
given part. There is no condition on how many arcs a part contains.

The condition is necessary for such a decomposition: for each ordered
cut, every strongly connected spanning part must supply an outgoing
arc, and the parts supply distinct arcs. The issue is whether this
necessary cut condition suffices in tournaments.

An equivalent formulation asks for $k$ arc-disjoint strongly connected
spanning subdigraphs without requiring them to exhaust $A$. Any remaining
arcs can be assigned to one part, because adding arcs preserves strong
connectivity. The assertion concerns an arc partition, rather than a
partition of vertices into strongly connected subtournaments.
\subsection{Short English statement}
Does every k-arc-strong tournament split its arcs into k strongly connected spanning subdigraphs?
\subsection{Research attempt: prime cyclic tournaments and the regular boundary}
\paragraph{Scope and literature check.}
GPT-6 Astra Ultra, 2 October 2026. Each part must use every vertex
and be strongly connected; vertices are not being partitioned.
Bang-Jensen and Yeo [S3] prove the two-arc-strong tournament case
and additional sufficient conditions for larger $k$. Those results
supply context without establishing the full universal cut-condition
claim. Here the attempt proves an explicit infinite family and
identifies exactly what the conjecture demands at the smallest
possible regular order.

\paragraph{1. Necessary cuts and a useful completion observation.}
Every strongly connected spanning digraph has an arc leaving each
nonempty proper vertex set. Thus $k$ arc-disjoint spanning strong
subdigraphs require at least $k$ arcs on each such cut. Singleton
sets and their complements in particular give both minimum indegree
and minimum outdegree at least $k$. In a tournament these neighbour
sets are disjoint, so the order is at least $2k+1$.

It is enough to construct the $k$ strong subdigraphs without covering
all arcs. Assign every unused arc to one chosen part; adding arcs
cannot destroy paths already present. This supplies an exact arc
partition with no further work. For $k=1$ the entire strongly
connected tournament is itself the required part.

\paragraph{2. At order $2k+1$, the target is Hamilton decomposition.}
Let $T$ be a regular tournament on $n=2k+1$ vertices, so every
outdegree is $k$. For a subset $X$ of size $s$, summing the
outdegrees in $X$ and subtracting its internal arcs gives
\[
 |A(X,V\setminus X)|=ks-\binom{s}{2}
                      =\frac{s(n-s)}2.
\]
For $1\le s\le n-1$, this concave quadratic is minimized at an
endpoint, where its value is $k$. Hence every such regular tournament
is $k$-arc-strong, independently of how its arcs are arranged.

There are exactly $nk$ arcs. Any spanning strong subdigraph has
at least $n$ arcs because every vertex must have an outgoing arc.
Therefore a partition into $k$ spanning strong parts forces exactly
$n$ arcs in each. Each vertex of a part then has exactly one outgoing
and, by the same count, one incoming arc. Such a digraph is a union
of directed cycles, and strong connectivity forces it to be one
Hamilton cycle. Conversely a Hamilton decomposition provides the
required partition. Thus, at this regular boundary, the proposed
statement is precisely a Hamilton decomposition assertion, rather
than merely a loose packing of arbitrary strong subgraphs.

\paragraph{3. A complete explicit family.}
Let $p=2k+1$ be an odd prime and use the cyclic tournament on
$\mathbb Z/p\mathbb Z$ whose arcs have steps $1,\ldots,k$.
For each $s$ in that range, put all arcs $x\to x+s$ in part $s$.
As $p$ is prime and $s\ne0$, successive additions of $s$ visit
every residue before returning to the start. This part is a
Hamilton cycle. The parts are disjoint because an ordered arc has
a unique step, and together they are all arcs of the tournament.
The cut calculation proves this tournament is $k$-arc-strong, so
these are explicit certificates under the original hypothesis.

The primality restriction belongs to this construction, not to the
conjecture. On nine vertices the same cyclic tournament has steps
one through four, but the step-three factor consists of three
separate directed triangles. It is not strong. This demonstrates
that decomposing into fixed-step factors need not work at composite
orders; it does not show that some other strong decomposition fails.
For a step $s$ modulo $n$, the exact number of cycles is
$\gcd(n,s)$, since each orbit has length $n/\gcd(n,s)$.

\paragraph{Verification and first remaining gap.}
The script [S9] verifies all cuts and the explicit partitions
for primes $3,5,7,11,13$, and checks the three components of the
step-three factor at order nine. The modular proof holds for every
odd prime. In general tournaments there is no translation-invariant
arc labelling to assign factors, and even in regular tournaments
arbitrary cycle factors may have several components. Converting
those factors to strong spanning parts while preserving disjointness
is the first gap in this approach.

\paragraph{Final status.}
UNRESOLVED generally. The prime cyclic family and the exact regular
boundary equivalence are proved; the known $k=2$ result is attributed
to [S3], not claimed as new work.
\subsection{Sources}
\begin{itemize}
\item[S1] ULAM.AI and the UnsolvedMath Contributors, supplied problems.json, record 3280 (OPG-47651), OpenGarden collection. Catalogue identity and source transcription; CC BY 4.0.
\url{https://huggingface.co/datasets/ulamai/UnsolvedMath}.
\item[S2] Open Problem Garden, Decomposing k-arc-strong tournament into k spanning strong digraphs. Original problem and discussion; the mathematical formulation below is expanded from the supplied source record.
\url{http://www.openproblemgarden.org/op/decomposing_k_arc_strong_tournament_into_k_spanning_strong_digraphs}.
\item[S3] J\o rgen Bang-Jensen and Anders Yeo, \emph{Decomposing $k$-arc-strong tournaments into strong spanning subdigraphs}, Combinatorica 24 (2004), 331--349. \url{https://doi.org/10.1007/s00493-004-0021-z}. Author-hosted manuscript: \url{https://pure.royalholloway.ac.uk/files/889922/paper_yeo_2.pdf}.
\item[S9] GPT-6 Astra Ultra, exact finite checks accompanying this attempt (2 October 2026). Repository script, Python standard library only. \texttt{scripts/check\_attempt\_batch03\_digraphs\_2026\_10\_02.py}.
\end{itemize}
\end{document}
